% =====================================================================
%  1.3 Rotuttrykk og brøkeksponenter — teori og oppgaver
% =====================================================================

\del{1}{Kvadratrot og n-te rot}

\begin{ordliste}
kvadratrot & Tallet $\sqrt{a}$ er det tallet som er $\ge 0$ og som ganget med seg selv gir $a$. & \mittsprak\\
$n$-te rot & Tallet $\sqrt[n]{a}$ er det tallet som opphøyd i $n$ gir $a$. For partall $n$: tallet $\ge 0$. Eksempel: $\sqrt[3]{8} = 2$. & \mittsprak\\
rotuttrykk & Et uttrykk med rottegn, for eksempel $3\sqrt{2} + 1$. & \mittsprak\\
irrasjonalt tall & Et tall som ikke kan skrives som en brøk av heltall. $\sqrt{2}$ er irrasjonalt. & \mittsprak\\
rasjonal nevner & En nevner uten rottegn. Å gjøre nevneren rasjonal betyr å fjerne roten fra nevneren. & \mittsprak\\
\end{ordliste}

\begin{mustkunne}
\begin{minipage}[t]{0.47\linewidth}
\textbf{Definisjon.} $\sqrt{9} = 3$, fordi $3 \ge 0$ og ${3 \cdot 3 = 9}$. Partallsrøtter av negative tall finnes ikke. Oddetallsrøtter finnes: $\sqrt[3]{-8} = -2$.\par
\textbf{Hvorfor produktregelen virker:}
\[ (\sqrt{a}\cdot\sqrt{b})^2 = (\sqrt{a})^2 (\sqrt{b})^2 = ab \]
Og $\sqrt{a}\cdot\sqrt{b} \ge 0$. Da er det tallet $\ge 0$ som gir $ab$ når det kvadreres, altså $\sqrt{ab}$.
\end{minipage}\hfill
\begin{minipage}[t]{0.49\linewidth}
\textbf{Regler} ($a, b \ge 0$)\par
\renewcommand{\arraystretch}{1.6}
\begin{tabular}[t]{@{}l l@{}}
$\sqrt{ab} = \sqrt{a}\cdot\sqrt{b}$ & Produkt\\
$\sqrt{\dfrac{a}{b}} = \dfrac{\sqrt{a}}{\sqrt{b}}$ & Brøk ($b \neq 0$)\\
$\sqrt{a}\cdot\sqrt{a} = a$ & Rot ganger rot\\
\end{tabular}
\end{minipage}
\end{mustkunne}

\begin{eksempel}[1 — forenkle og trekke sammen]
Forenkle $\sqrt{48} - \sqrt{12} + \sqrt{27}$.

\begin{stegtabell}
$= \sqrt{16\cdot 3} - \sqrt{4\cdot 3} + \sqrt{9\cdot 3}$ & Finn største kvadrattall som er faktor i hvert tall.\\
$= 4\sqrt{3} - 2\sqrt{3} + 3\sqrt{3}$ & Produktregelen: $\sqrt{16\cdot3} = \sqrt{16}\cdot\sqrt{3}$.\\
$= (4 - 2 + 3)\sqrt{3} = 5\sqrt{3}$ & Like rotledd trekkes sammen, som $4x - 2x + 3x$.\\
\end{stegtabell}
\end{eksempel}

\begin{eksempel}[2 — rasjonal nevner]
Gjør nevneren rasjonal: $\dfrac{10}{\sqrt{5}}$.

\begin{stegtabell}
$= \dfrac{10\cdot\sqrt{5}}{\sqrt{5}\cdot\sqrt{5}} = \dfrac{10\sqrt{5}}{5} = 2\sqrt{5}$ & Utvid med $\sqrt{5}$. Nevneren blir 5. Forkort.\\
Kontroll: & $2\sqrt{5}\cdot\sqrt{5} = 2\cdot 5 = 10$. Da er $\frac{10}{\sqrt{5}} = 2\sqrt{5}$.\\
\end{stegtabell}
\end{eksempel}

\begin{vanligfeil}
\begin{itemize}[leftmargin=1.2em,itemsep=1pt,topsep=0pt]
  \item $\sqrt{a+b} = \sqrt{a} + \sqrt{b}$ er feil. Sjekk: $\sqrt{9+16} = 5$, men $3 + 4 = 7$.
  \item $\sqrt{8} + \sqrt{2} = \sqrt{10}$ er feil. Forenkle først: $2\sqrt{2} + \sqrt{2} = 3\sqrt{2}$.
  \item $2\sqrt{3} = \sqrt{6}$ er feil. Tallet 2 går under roten som $2^2$: $2\sqrt{3} = \sqrt{4\cdot 3} = \sqrt{12}$.
\end{itemize}
\end{vanligfeil}

\nivaahode{1}

\oppg{1.1}{1}{\uh}
Regn ut.
\begin{tasks}(4)
  \task $\sqrt{81}$
  \task $\sqrt[3]{27}$
  \task $\sqrt[3]{-125}$
  \task $\sqrt{\dfrac{4}{25}}$
\end{tasks}

\oppg{1.2}{1}{\uh}
Forenkle.
\begin{tasks}(4)
  \task $\sqrt{28}$
  \task $\sqrt{50}$
  \task $\sqrt{45}$
  \task $\sqrt{72}$
\end{tasks}

\nivaahode{2}

\oppg{1.3}{2}{\uh}
Forenkle.
\begin{tasks}(2)
  \task $\sqrt{18} + \sqrt{50} - \sqrt{8}$
  \task $\sqrt{63} + \sqrt{28}$
  \task $3\sqrt{20} - \sqrt{45}$
  \task $\sqrt{2}\cdot\sqrt{18}$
\end{tasks}

\oppg{1.4}{2}{\uh}
Gjør nevneren rasjonal. Forenkle svaret.
\begin{tasks}(4)
  \task $\dfrac{6}{\sqrt{3}} + \sqrt{27}$
  \task $\dfrac{3}{2\sqrt{3}}$
  \task $\dfrac{4}{\sqrt{8}}$
  \task $\dfrac{10}{\sqrt{20}}$
\end{tasks}

\oppg{1.5}{2}{\uh}
Finn \textbf{den første feilen}. Regn riktig.
\[ \sqrt{2}\left(\sqrt{8}+\sqrt{18}\right) \ \to\ \sqrt{2}\cdot\sqrt{26} \ \to\ \sqrt{52} \ \to\ 2\sqrt{13} \]

\nivaahode{3}

\oppg{1.6}{3}{\uh}
Er påstanden alltid, noen ganger eller aldri sann? Begrunn svaret.
\begin{tasks}(1)
  \task $\sqrt{x^2} = x$ for et vilkårlig tall $x$
  \task $\sqrt{a+b} = \sqrt{a} + \sqrt{b}$ for $a, b \ge 0$
  \task $\sqrt{a}+\sqrt{a} = \sqrt{2a}$ for $a \ge 0$
\end{tasks}

\oppg{1.7}{3}{\uh}
Anta at $a, b \ge 0$.
\begin{tasks}(2)
  \task Vis at $\sqrt{a^2 b} = a\sqrt{b}$.
  \task Forenkle $\sqrt{75a^2}$.
\end{tasks}

\begin{sjekk}
\begin{tasks}(2)
  \task Forenkle $\sqrt{32}$.
  \task Forenkle $\sqrt{75} - \sqrt{12}$.
  \task! Er $2\sqrt{3} = \sqrt{6}$? Forklar.
\end{tasks}
\end{sjekk}

% ---------------------------------------------------------------------

\del{2}{Brøk som eksponent}

\begin{ordliste}
brøkeksponent & En eksponent som er en brøk, for eksempel $9^{\tfrac{1}{2}}$ eller $8^{\tfrac{2}{3}}$. & \mittsprak\\
positivt grunntall & Brøkeksponenter bruker vi bare når grunntallet $a$ er større enn 0. & \mittsprak\\
omskrive & Å skrive det samme tallet på en annen form, for eksempel en rot som en potens. & \mittsprak\\
\end{ordliste}

\begin{mustkunne}
\textbf{Hvorfor $a^{\tfrac12} = \sqrt{a}$?} Potensregelen $(a^m)^n = a^{m\cdot n}$ skal gjelde også for brøker. Da må
\[ \left(a^{\tfrac12}\right)^2 = a^{\tfrac12 \cdot 2} = a^1 = a \]
Altså er $a^{\tfrac12}$ et positivt tall som ganget med seg selv gir $a$. Det er definisjonen av $\sqrt{a}$.
Vi definerer brøkeksponenter bare for $a > 0$:
\[ a^{\tfrac{1}{n}} = \sqrt[n]{a} \qquad\qquad a^{\tfrac{m}{n}} = \left(\sqrt[n]{a}\right)^m = \sqrt[n]{a^m} \qquad\qquad a^{-\tfrac{m}{n}} = \frac{1}{a^{\tfrac{m}{n}}} \]
Nevneren i eksponenten er rotens tall. Telleren er potensen.
\end{mustkunne}

\begin{eksempel}[1 — regn ut $16^{\tfrac34}$ på to måter]
\begin{stegtabell}
$16^{\tfrac34} = \left(\sqrt[4]{16}\right)^3 = 2^3 = 8$ & Rot først: $2^4 = 16$. Tallene er små.\\
$16^{\tfrac34} = \sqrt[4]{16^3} = \sqrt[4]{4096} = 8$ & Potens først: vi må kjenne igjen at $8^4 = 4096$.\\
\end{stegtabell}
Begge gir 8. Uten hjelpemidler er den første måten lettest: ta roten først.
\end{eksempel}

\begin{eksempel}[2 — skriv som én potens]
Skriv $\dfrac{x\sqrt{x}}{\sqrt[3]{x}}$ som én potens av $x$ ($x > 0$).

\begin{stegtabell}
$= \dfrac{x^1 \cdot x^{\tfrac12}}{x^{\tfrac13}}$ & Skriv hver rot som en potens.\\
$= x^{1 + \tfrac12 - \tfrac13}$ & Produkt: legg sammen. Kvotient: trekk fra.\\
$= x^{\tfrac{6}{6}+\tfrac{3}{6}-\tfrac{2}{6}} = x^{\tfrac76}$ & Felles nevner 6.\\
Kontroll med $x = 64$: & $\frac{64\cdot 8}{4} = 128$ og $64^{\tfrac76} = 2^7 = 128$. Det stemmer.\\
\end{stegtabell}
\end{eksempel}

\begin{vanligfeil}
\begin{itemize}[leftmargin=1.2em,itemsep=1pt,topsep=0pt]
  \item $a^{\tfrac12} = \frac{a}{2}$ er feil. Eksponenten $\frac12$ betyr kvadratrot: $a^{\tfrac12} = \sqrt{a}$.
  \item $8^{\tfrac23} = \frac{8^2}{3}$ er feil. Nevneren 3 er rotens tall: $8^{\tfrac23} = (\sqrt[3]{8})^2 = 4$.
  \item $a^{-\tfrac12} = -\sqrt{a}$ er feil. Negativ eksponent gir brøk, som i 1.1: $a^{-\tfrac12} = \frac{1}{\sqrt{a}}$.
\end{itemize}
\end{vanligfeil}

\nivaahode{1}

\oppg{2.1}{1}{\uh}
Skriv som en potens med brøkeksponent ($x > 0$).
\begin{tasks}(4)
  \task $\sqrt{7}$
  \task $\sqrt[3]{5}$
  \task $\sqrt{x^3}$
  \task $\left(\sqrt[3]{x}\right)^2$
\end{tasks}

\oppg{2.2}{1}{\uh}
Regn ut.
\begin{tasks}(4)
  \task $25^{\tfrac12}$
  \task $27^{\tfrac13}$
  \task $4^{\tfrac32}$
  \task $64^{\tfrac23}$
\end{tasks}

\nivaahode{2}

\oppg{2.3}{2}{\uh}
Regn ut.
\begin{tasks}(4)
  \task $27^{-\tfrac13}$
  \task $81^{-\tfrac34}$
  \task $\left(\dfrac{4}{9}\right)^{-\tfrac12}$
  \task $32^{\tfrac25}$
\end{tasks}

\oppg{2.4}{2}{\uh}
Skriv som én potens av $x$ ($x > 0$).
\begin{tasks}(4)
  \task $x\cdot\sqrt{x}$
  \task $\sqrt{x}\cdot\sqrt[3]{x}$
  \task $\dfrac{x^2}{\sqrt[4]{x}}$
  \task $\dfrac{\sqrt[3]{x^2}}{\sqrt{x}}$
\end{tasks}

\oppg{2.5}{2}{\uh}
Finn \textbf{den første feilen}. Regn riktig.
\[ 8^{-\tfrac23} \ \to\ \left(\sqrt[3]{8}\right)^{-2} \ \to\ 2^{-2} \ \to\ -4 \]

\oppg{2.6}{2}{\mh}
Keplers tredje lov sier at omløpstiden er $T = a^{\tfrac32}$. Her er $T$ i år og $a$ er middelavstanden til sola i astronomiske enheter.
\begin{tasks}(1)
  \task Mars har $a \approx 1,52$. Finn omløpstiden. Svar med to gjeldende siffer.
  \task Et objekt har $T = 8$ år. Finn $a$ eksakt.
\end{tasks}

\nivaahode{3}

\oppg{2.7}{3}{\uh}
\begin{tasks}(1)
  \task Løs likningen $x^{\tfrac32} = 27$.
  \task For hvilke positive heltall $n$ er $n^{\tfrac32}$ et heltall? Begrunn.
\end{tasks}

\oppg{2.8}{3}{\uh}
La $a > 0$. Er påstanden alltid, noen ganger eller aldri sann? Begrunn svaret.
\begin{tasks}(1)
  \task $a^{-\tfrac12} = -\sqrt{a}$
  \task $a^{\tfrac12} = \dfrac{a}{2}$
  \task $a^{\tfrac12} < a$
\end{tasks}

\begin{sjekk}
\begin{tasks}(2)
  \task Regn ut $9^{\tfrac32}$.
  \task Regn ut $125^{-\tfrac13}$.
  \task! Forklar hvorfor $a^{\tfrac12} = \sqrt{a}$. Bruk regelen $(a^m)^n = a^{mn}$.
\end{tasks}
\end{sjekk}
